\section{Reduced spectral space and parameter hierarchy}
\label{sec:retained}

Choose the parameters in the following order. First fix any constant
$C_{\rm pkt}>0$, which determines the logarithmic packet reserve
\eqref{eq:packet-bandwidth-reserve} before any stationary-buffer or symbol
order is chosen. By Proposition~\ref{prop:trivial-outer-twist}, no arithmetic
range parameter is required.
% Canonical parameter guard retained for the proof-verification contract:
% 0<\kappa<\frac1{10}.

Fix once and for all the finite derivative orders required for two curvature
derivatives and two primitive mixed derivatives. Let $C_{\rm fr}$ be the
freezing exponent in Lemma~\ref{lem:block-freezing}, and choose a fixed height
partition exponent
\begin{equation}
D>C_{\rm fr}+5,
\qquad H:=T/L^D.
\label{eq:height-block-choice}
\end{equation}
Let $C_{\rm mean}$ dominate the fixed local-source/refinement multiplicity, and
let $C_{\rm reg}$ dominate the per-shell logarithmic exponents of the fixed
lag-profile, freezing, $f'$, and stationary/Stirling correction classes. It
does not include the nonstationary integration order or the high-sideband
Fourier order, whose exponents $C_N$ and $C_{\rm sb}$ are introduced only after
those orders are chosen below. Let $C_{\rm diag}$ dominate all fixed packet,
curvature, and shell losses multiplying the arbitrary-log term in
\eqref{eq:packet-HLZ-error}, and fix
\begin{equation}
A_{\rm diag}>C_{\rm diag}+10.
\label{eq:Adiag-choice}
\end{equation}
Choose
\[
\sigma>2,
\qquad \epsilon_L=L^{-\sigma}.
\]
The direct high-sideband estimate below is obtained with the explicit shell
power $d_j^3$; no additional shell-dimension exponent is assumed.
Next fix $M_*>2C_{\rm reg}+5$. We use the same numerical target as a
logarithmic saving exponent for the stable-core nonstationary leakage and, after
increasing it if necessary, as a power-saving target for the horizontal
connectors. Choose the integration-by-parts order $N$ so that
\[
2N-1>M_*+5,
\]
and then choose the stationary buffer exponent $B$ so that
\[
B(2N-1)>2C_N+M_*.
\]
Independently choose one finite high-sideband order
\[
M_{\rm sb}>5,
\]
and let $C_{\rm sb}$ be the corresponding fixed logarithmic exponent in
\eqref{eq:high-sideband-one-shell}. Finally choose $A_0$ so large that
\begin{equation}
A_0>\max\left\{
B+3,
\sigma+3,
C_{\rm mean}+3,
C_{\rm reg}+3,
\frac{C_{\rm sb}+\sigma M_{\rm sb}+3}
{M_{\rm sb}-2}
\right\}.
\label{eq:acyclic-A0-choice}
\end{equation}
With $W_T,X_T$ defined by \eqref{eq:stable-log-core}, these derivative orders
and exponents are fixed once and for all before taking the limit $T\to\infty$.

\begin{lemma}[Global and local mean constraints]
\label{lem:mean-zero-space}
Let $g_\alpha(u;v)$, $0\le u\le W_T$, denote the output of the common flat
analysis map immediately before the global primitive in
\eqref{eq:global-reference-Gram}, with $\alpha$ collecting the retained
arithmetic and fixed local-source labels. Define the global mean map
\begin{equation}
\mathcal M_T^{\rm glob}v
:=\left(\int_0^{W_T}g_\alpha(u;v)\,du\right)_\alpha.
\label{eq:global-mean-map}
\end{equation}
After the logarithmic unitary write the corresponding output as
$f_\alpha(x;v)$, $1\le x\le X_T$. On
$I_j=[J_j,J_j+1]$ set
\begin{equation}
q_{j,\alpha}(t;v)
:=\psi_j(J_j+t)(J_j+t)^{-1/2}
 f_\alpha(J_j+t;v),
\qquad 0\le t\le1,
\label{eq:local-primitive-input}
\end{equation}
and define the local mean map
\begin{equation}
\mathcal M_T^{\rm loc}v
:=\left(\int_0^1q_{j,\alpha}(t;v)\,dt\right)_{j,\alpha}.
\label{eq:local-mean-map}
\end{equation}
Finally put
\begin{equation}
\boxed{
V_{\rm mean}:=\ker\mathcal M_T^{\rm glob}
\cap\ker\mathcal M_T^{\rm loc}.
}
\label{eq:mean-map}
\end{equation}
Then
\begin{equation}
\boxed{
\codim V_{\rm mean}
\ll (1+X_T)L^{C_{\rm mean}}
=o(N_T).
}
\label{eq:mean-codim}
\end{equation}
For every $v\in V_{\rm mean}$, the global primitive
\[
R_\alpha(u;v):=\int_0^u g_\alpha(s;v)\,ds
\]
satisfies $R_\alpha(0;v)=R_\alpha(W_T;v)=0$, while every local primitive
\[
Q_{j,\alpha}(t;v):=\int_0^tq_{j,\alpha}(s;v)\,ds
\]
satisfies $Q_{j,\alpha}(0;v)=Q_{j,\alpha}(1;v)=0$. Thus the global model
primitive identity and the shell-local remainder identity hold on the same
retained subspace.
\end{lemma}

\begin{proof}
Both mean maps are linear. By Proposition~\ref{prop:trivial-outer-twist}, the
invariant source has a single arithmetic fibre; the remaining fixed local-source
labels contribute at most $L^{C_{\rm mean}}$ coordinates. There are $O(X_T)$
retained unit shells, so the target of $\mathcal M_T^{\rm loc}$ has
$O(X_TL^{C_{\rm mean}})$ coordinates. Hence the codimension of the
intersection of their kernels is bounded by the sum of these dimensions. The
choice $A_0>C_{\rm mean}+3$ in \eqref{eq:acyclic-A0-choice} makes this
$o(TL)=o(N_T)$. The endpoint identities
for $R_\alpha$ and $Q_{j,\alpha}$ are exactly the two kernel conditions in the
definition of $V_{\rm mean}$.
\end{proof}

All edge-rank, shell-row, and regular-remainder estimates above were proved on
the full admissible near-lattice packet family. Passing to the
stationary-node-sized subfamily selected in Lemma~\ref{lem:generic-affine} is a
coordinate compression. Rank and absolute Hilbert--Schmidt norm cannot
increase, and the uniform output-fibre estimates used for the relative bounds
may be rerun on the compressed packet family with the same asymptotics.

It remains only to impose the horizontal endpoint restriction. Let
$P_{\rm end}$ be the coordinate projection which removes the packet centres in
the two collars of Lemma~\ref{lem:horizontal-decay}, and set
\begin{equation}
V_{\rm end}:=\Ran P_{\rm end}.
\label{eq:end-space}
\end{equation}
Then
\begin{equation}
\boxed{\codim V_{\rm end}=O(T)=o(N_T).}
\label{eq:end-codim}
\end{equation}
This is a genuine packet-coordinate restriction, so
Lemma~\ref{lem:coordinate-restriction} applies after intersection with
$V_{\rm end}$. The stationary buffer rows are not deleted. Their contribution is accounted
for by the $R_j$ enlargement in the edge-rank estimate and by the complementary
nonstationary remainder bounds.

Define the retained subspace by
\begin{equation}
\boxed{
\Pret=V_{\rm core}\cap V_{\rm mean}\cap V_{\rm end}.
}
\label{eq:retained-space}
\end{equation}
By Lemma~\ref{lem:stable-core-frame}, \eqref{eq:mean-codim}, and
\eqref{eq:end-codim},
\begin{equation}
\boxed{\codim\Pret=o(N_T).}
\label{eq:retained-codim}
\end{equation}
The frozen model requires no further spectral deletion by
Proposition~\ref{prop:model-floor}. Stationary shell rows, the smooth height
partition, the outer excess, and the low-sideband ranges are not additional
subspace restrictions.
